\section{Which interval to report}\label{sec:q1}\label{sec:impl}

Consider two groups of firms joined by a small number of workers who switch
between them. Many wages can pin down relative pay within each group, while
those few movers determine the difference in pay between the groups. Firm
variance depends on the \emph{square} of that difference. Even if the estimated
difference is nearly normal, its square can be far from normal. A symmetric
interval based on the square root of \eqref{eq:Vhat} can then be too short or
even unavailable. KSS treat this case with a separate interval, which they index by
$q=1$. (Their $q$ counts dominant directions. It is unrelated to the sample variance
$q_g$.) The variance inputs of Section~\ref{sec:estimator} supply everything that
interval needs.

\fevc{} reports one interval per component. In a well-connected design a positive
variance estimate gives the familiar interval
$\hat\theta\pm1.96\sqrt{\widehat{\V}[\hat\theta]}$ at the 95\% level.
When a single direction accounts for more than 10\% of the estimated sampling
variance, the program reports the $q=1$ interval if it is available. It also
uses that interval when the normal variance estimate is nonpositive.
For the variance of the sum the nonzero eigenvalues of the target are all
equal, so no direction can dominate, and the default keeps the normal interval when
available. The choice is made component by component: the same mobility network can
support a normal interval for one component and call for the $q=1$ interval for
another.

On large, well-connected panels the normal interval is the rule. On three Veneto
panels of 2.3 to 5.2 million person-years, under both deletion modes, the leading
direction's share was between 0.03\% and 8.1\% and every component received the
normal interval. In the simulations the switch is what keeps coverage when a
bottleneck is present (Section~\ref{sec:impl-mc}).

\subsection{One dominant direction}\label{sec:q1int}
For each component, take the eigenvalue $\lambda_1$ of $Av=\lambda S_{xx}v$
that is largest in absolute value. Normalise the corresponding observation
weights $w_1=Xv_1$ to unit length. The scalar score $\hat b_1=w_1'y$
estimates $b_1=w_1'\mu$. With exact weights and without estimated controls,
KSS's leave-out variance estimate for this score is
$V_b^{\rm lo}=\sum_i w_{i1}^2\hat\sigma_i^2$. Removing the direction from
the target gives the identity
\[
  \hat\theta=\lambda_1(\hat b_1^2-V_b^{\rm lo})+\hat\theta_{\rm r}.
\]
The remainder $\hat\theta_{\rm r}$ is another leave-out estimate. Its weights
$C^{\rm r}$ follow from \eqref{eq:ustat} with $B-\lambda_1w_1w_1'$ in place of
$B$. Write $m^{\rm r}=C^{\rm r}\mu$ and $\hat m^{\rm r}=C^{\rm r}y$.
The covariance matrix of the score and the remainder has entries
\begin{equation}\label{eq:q1cov}
 \Sigma_{bb}=\sum_i w_{i1}^2\sigma_i^2,\qquad
 \Sigma_{b{\rm r}}=2\sum_i w_{i1}m_i^{\rm r}\sigma_i^2,\qquad
 \Sigma_{\rm rr}=\V[\hat\theta_{\rm r}],
\end{equation}
where the last variance is \eqref{eq:V} for $C^{\rm r}$. The package estimates
the entries with the same variance inputs used for the normal interval:
\begin{equation}\label{eq:q1est}
 V_b=\hat\Sigma_{bb}=\max\!\left\{V_b^{\rm lo},\,\sum_i w_{i1}^2\st_i^2\right\},
 \qquad \hat\Sigma_{b{\rm r}}=2\sum_i w_{i1}\hat m_i^{\rm r}\st_i^2,\qquad
 \hat\Sigma_{\rm rr}=\widehat{\V}[\hat\theta_{\rm r}].
\end{equation}
Here $\widehat{\V}[\hat\theta_{\rm r}]$ includes the pair adjustment with the
remainder's weights.

\paragraph{Why the larger of two numbers.}
The two candidates for $V_b$ protect against different errors. The leave-out
estimate is KSS's. It is unbiased and needs no model for anyone, because
$\Sigma_{bb}$ is linear in the error variances. But with a bottleneck the score rests
on the few rows that link the groups, and the leave-out products of a few rows are
noisy: their weighted sum can come out far too small, or negative. The fitted estimate
is positive and stable, but only as good as the model on those few rows. Taking the
larger guards against a downward error in either. This is why a model that understates
the variance of the linking movers shortens the normal standard error but not the
$q=1$ interval (Section~\ref{kss:sec:synth}). The price is, as a rule, a longer
interval. The leave-out estimate $V_b^{\rm lo}$ still enters the decomposition of the
point estimate, which is KSS's.

\paragraph{The interval.}
The interval takes the smallest and largest values of $\lambda_1 b^2+\theta_{\rm r}$
over a confidence ellipse for $(b,\theta_{\rm r})$ centred at
$(\hat b_1,\hat\theta_{\rm r})$. It thus allows for the curvature of the
square. With $\hat\rho$ the estimated correlation in \eqref{eq:q1est}, the curvature is
\begin{equation}\label{eq:kappa}
 \hat\kappa=\frac{2|\lambda_1|V_b}
 {\sqrt{\hat\Sigma_{\rm rr}(1-\hat\rho^2)}} .
\end{equation}
Writing $u=(b-\hat b_1,\theta_{\rm r}-\hat\theta_{\rm r})'$, the ellipse
is $u'\hat\Sigma^{-1}u\le z_{\alpha,\hat\kappa}^{\,2}$.
The Andrews--Mikusheva (2016) radius $z_{\alpha,\kappa}$ is the
$1-\alpha$ quantile of
$\sqrt{Z_1^2+(|Z_2|+\kappa^{-1})^2}-\kappa^{-1}$ for independent
standard normal $Z_1,Z_2$, interpreted by continuity at zero curvature.
The package simulates the critical value with 100,000 draws. Its 95\% radius rises
from 1.96 at zero curvature towards 2.45 at high curvature, so without a dominant
direction the two intervals nearly coincide. A positive-definite covariance matrix is
necessary. The reported endpoints can be asymmetric, and their distance from the point
estimate cannot be summarised by a single standard error.

\subsection{The choice made by the command}
When the normal variance estimate is positive, \fevc{} compares 0.10 with the
\emph{lead share}
\[
 \frac{4\lambda_1^2\max(\hat b_1^2-V_b,0)V_b+2\lambda_1^2V_b^2}
      {\widehat{\V}[\hat\theta]} ,
\]
which the output table displays for every component but the variance of the sum. The
numerator estimates the
sampling variance attributable to the leading direction, including the part that
depends on the size of its score. Both numerator and denominator are estimated, so the
ratio can exceed one. KSS suggest a heuristic threshold instead: treat as dominant the
directions whose squared eigenvalue is at least one tenth of the sum of squared
eigenvalues. That statistic uses only the design. The two answer different questions,
even though both cutoffs are 10\%.

The variance share was preferred for three reasons.
\begin{itemize}
\item \emph{It measures what causes the problem.} The departure from normality comes
from the leading direction's part of the sampling variance. That part depends on the
error variances along the direction and on the size of its score, not only on the
design: a weakly connected direction can dominate the eigenvalues and still carry
little of the variance.
\item \emph{It reports the $q=1$ interval only where it is needed.} A $q=1$ component
has no standard error and leaves \code{e(V)} unposted. On the stored draws of the
$q=1$ simulations (27 designs, 1,000 draws each) both rules kept coverage, with a
lowest coverage of 0.940 for KSS's heuristic and 0.937 for the variance share. But
KSS's heuristic chose $q=1$ for 78 of the 81 design--component pairs, including the firm variance
and the covariance of a design without a bottleneck, whose eigenvalue shares were 0.11
and 0.12. The variance share chose it in 7\% of the draws of designs without a
bottleneck. This comparison was made after the results of those simulations were
known, so it is development evidence.
\item \emph{It is free.} The variance share uses only quantities that the $q=1$
calculation already produces. The eigenvalue share needs a separate randomised trace
of the squared target kernel (KSS's code uses 10,000 draws).
\end{itemize}
Its drawbacks are that it is a ratio of estimates, so it is noisy, and that it
inherits the variance model's errors.

\paragraph{What the user sees.}
The output shows, for every component, the estimate, the standard error, the interval,
which interval it is, and the lead share (none for the variance of the sum). For a $q=1$ component the standard-error
column is empty. \code{e(inference\_intervals)} holds both intervals of every component
and the normal-form standard error, and \code{estat inference} explains the choice.
When all four components have normal intervals and their joint covariance is
available, \code{e(V)} gives that covariance. Its fourth row and column respect the
identity between the variance of the sum and the three base components.
\code{inference(highrank)} requests the normal form for every component and still
reports the lead shares when the search for the directions succeeds. \code{inference(q1)} requests the $q=1$ form wherever it is
available. Under the default a request that cannot support inference returns point
estimates with a reason, and \code{inference(none)} requests point estimates alone.

\paragraph{Two or more dominant directions.}
The command does not implement an interval for more than one dominant direction.
\code{estat inference} reports the two largest eigenvalues and an effective rank for
each component and warns when the rank is low. A second eigenvalue close to the first
is the sign that neither interval's conditions hold.

\paragraph{What supports the choice.}
The coverage result for a prespecified $q=1$ interval requires a diffuse
remainder, a leading score supported by a growing effective number of
observations, accurate variance inputs in the relevant directions, and the
conditions for joint normality. For the larger-of-two rule in \eqref{eq:q1est},
the result also bounds the correlation between score and remainder
(Corollary~\ref{cor:q1}). The 10\% switching rule itself was chosen on development
simulations and confirmed on fresh designs. Its coverage is an empirical
finding (Section~\ref{sec:impl-mc}), and neither KSS nor this note has a theorem for
a data-dependent choice of interval.

\section{Controls, centring and computation}\label{sec:controls}

\subsection{Which estimator the control adjustment describes}
Year effects illustrate why estimated controls matter. Removing the true year
effects would leave repeated observations of the same worker at the same firm
with a common expected wage. Removing \emph{estimated} year effects also puts
the estimation error of each year coefficient into every adjusted wage for
that year. Treating those wages as independent can materially understate
uncertainty. \fevc{} adds a first-order term for the estimation of the control
coefficients, and the term does its job. In the dense Gaussian example of
Appendix~\ref{app:ctlnum}, holding the control estimates fixed gives only 62\% of the
true \emph{variance} of the estimated variance of the sum under match deletion with
one error variance per cell. The adjusted calculation gives 99.8\%.

The standard error describes the \emph{fixed-offset} estimator. This estimator
first fits the controls, removes their fitted contribution to wages, and then
performs leave-out estimation on the adjusted outcome. Its uncertainty includes
estimation of the controls to first order. The default joint point estimate
instead re-estimates the control coefficients inside each leave-out fit. That
adds a term which falls outside the reported standard error. In the dense check, under
observation deletion, the joint estimate's variance is 0.995 to 1.035 times the
fixed-offset one. The term can
be material when the control index is large or closely aligned with a difference
between groups of firms (Section~\ref{sec:impl-mc}). Under match deletion,
\code{nuisance(fixedoffset)} gives the point estimator to which the adjustment
applies directly. Under observation deletion the supported inference route requires
joint nuisance estimation, and the native route does not accept frequency weights.

To see the origin of the correction, let $v$ be the error vector in the
outcome after subtracting the fitted control index, and $S=\V[v]$. Although
the original deletion units are independent, $S$ generally has off-diagonal
entries. Under Gaussian errors the exact variance of the fixed-offset
quadratic form is $4m'Sm+2\tr(CSCS)$. The program adjusts the linear
influences and applies the corresponding control adjustment to simulated
noise. The feasible formula uses residual products to estimate the
within-unit control terms. Its trace calculation omits the covariance between
a unit error and its within-unit control product. Also, a set still receives
one pooled variance even when its control-adjusted weights differ within the
cell. Error variances that move with a control inside a match, poorly identified
control coefficients, or a dominant direction aligned with the controls can
make these approximations consequential. Appendix~\ref{app:controls} gives
the derivation and its seven approximations. The dense check uses true variances
outside repeated cells. End-to-end evidence with fitted variances appears in
Section~\ref{sec:impl-mc}.

\subsection{Centring and the cost of inference}\label{sec:computation}
Under the default \code{centering(mean)}, the program centres the simulated
noise in the same way as the outcome. Under \code{centering(none)} it uses
the uncentred weights. Under \code{centering(corrected)} it treats the
increment for the centring correction as fixed when computing the standard
error. These choices concern the sampling variability of the econometric
estimate. The Monte Carlo standard error of the \emph{point estimate}, also
reported by \fevc{}, measures numerical variation from random projections and is a
separate quantity.

The most expensive-looking term in \eqref{eq:Vhat} sums over every pair of
rows. It is computed with a few linear solves instead. In the baseline
case draw independent $\zeta_i$ with mean zero and variance $\st_i^2$, apply
$C$ once, and sum the squared entries of $C\zeta$ with weights $\st_i^2$.
Conditional on the fitted variances,
\begin{equation}\label{eq:sim}
 \E_\zeta\!\left[\sum_i\st_i^2(C\zeta)_i^2\right]
      =\sum_{i,\ell}C_{i\ell}^2\st_i^2\st_\ell^2 .
\end{equation}
The calculation uses normal draws and reuses them across components, their
covariances and the remainder. With many diffuse directions, each draw
already averages noise across many rows, so few draws are needed; the default is
eight. Proposition~\ref{prop:sim} relates the precision of this average to the
concentration of the spectrum. On the six-year Veneto panel, where the method study
used 16 draws, other sets of draws or a run with 64 draws changed no standard error by
more than 0.6\%.

The native route reuses the point estimator's leverages and target weights,
with 200 random projections by default and the finite-projection correction
in the point estimate. The standard errors then add about 90 linear solves to the
roughly 600 to 800 of the point estimates: three plus four per noise draw for the
normal covariance (35 with the default eight draws), and about 55 for the leading
directions and the $q=1$ pieces. The default, \code{inference(highrank)} and
\code{inference(q1)} all take them, because
\code{inference(highrank)} also searches for the leading directions in order to report
the lead shares. The native and matrix-free Mata routes find the directions by a
restarted Lanczos iteration. Exact Mata uses a dense eigendecomposition, while Mata
with random projections runs its own projection stage for inference.

On three Veneto panels with 2.3 to 5.2 million person-years, no controls and 16
threads, the default added 24--30\% to the measured running time of point estimation.
With controls under the default \code{nuisance(joint)}, forming the fixed-offset
version of the units takes one further pass of two solves per point-estimate
projection, 400 at the default. This note has no measured running time for that case.
By the
solve counts the offset pass is then the larger part of the cost of inference, and
\code{nuisance(fixedoffset)} avoids it.
